\chapter{Problem Formulation and Experimental Framework}
\label{ch:problem}

\section{Formal Problem Definition}

Let $D = \{x_1, x_2, \dots, x_N\}$ denote a collection of $N$ classical music recordings drawn from the iPalpiti archive.
Each item $x_i$ represents a complete audio track.

An embedding model is defined as a function:

\[
f_\theta : \mathcal{X} \rightarrow \mathbb{R}^d
\]

where $\theta$ denotes the model parameters and $d$ the embedding dimensionality.
The embedding of item $x_i$ is:

\[
z_i = f_\theta(x_i)
\]

Similarity between two items is computed using a similarity function $s(\cdot,\cdot)$, such as cosine similarity:

\[
s(z_q, z_i) = \frac{z_q \cdot z_i}{\|z_q\|\|z_i\|}
\]

Higher similarity values indicate greater embedding-space proximity.

Given a query item $x_q \in D$, all items in $D \setminus \{x_q\}$ are ranked in descending order of similarity $s(z_q, z_i)$. 
Retrieval quality is evaluated with respect to a relevance function $r_q(x_i) \in \{0,1\}$ defined by proxy tasks.
Ranking quality is measured using standard IR metrics including NDCG@K, Precision@K, and Recall@K.



\section{Definition of Embedding Capacity}

The primary independent variable in this study is embedding-model capacity. 
Embedding capacity is operationalized primarily through model parameter count ($|\theta|$), which serves as the principal quantitative measure of model size. 
Embedding dimensionality ($d$) is reported as a secondary structural attribute.

Architectural class (e.g., CNN-based vs.\ Transformer-based) is treated as a separate categorical factor rather than as a direct measure of capacity. 
Accordingly, models are compared at discrete capacity levels, with architectural differences analyzed independently from raw parameter count. 
This operationalization allows capacity to be treated as a measurable experimental variable rather than a qualitative label (e.g., ``heavy'' vs.\ ``light'').

Capacity effects are examined within each architectural family via small, base, and large variants, enabling controlled scaling comparisons under shared structural principles.

This design enables empirical characterization of scaling behavior under constrained data conditions.

\section{Proxy Retrieval Tasks}

Relevance is defined via structured proxy tasks derived from metadata.
Let $g(x)$ denote a metadata-based labeling function.
Two proxy families are considered:

\paragraph{Sanity Proxy.}
Relevance is defined by shared composer identity:
\[
r_q(x_i) = 1 \quad \text{if} \quad g_{\text{composer}}(x_q) = g_{\text{composer}}(x_i)
\]

\paragraph{Primary Musical Proxy.}
Relevance is defined by shared character labels:
\[
r_q(x_i) = 1 \quad \text{if} \quad g_{\text{character}}(x_q) = g_{\text{character}}(x_i)
\]

For experimental tractability, relevance is treated as binary:
\[
r_q(x_i) \in \{0,1\}
\]
with $r_q(x_i)=0$ when labels do not match.

Although musical similarity is inherently continuous, metadata-derived labels provide a controlled and reproducible approximation suitable for systematic evaluation.
These proxy tasks serve as diagnostic instruments rather than end goals, enabling controlled analysis of embedding-space organization in constrained archival settings.


\section{Experimental Variables and Controls}

This study follows a hypothesis-driven experimental design.
Embedding-model capacity serves as the primary independent variable, while ranking-based retrieval metrics (NDCG@K, Precision@K, and Recall@K) constitute the dependent variables.
To ensure internal validity, the audio preprocessing pipeline, indexing mechanism, similarity function, and evaluation protocol are held constant across all experimental conditions.
Under this controlled setting, performance differences can be attributed primarily to variations in embedding capacity.


\section{Backend Pipeline Architecture}

The backend evaluation pipeline consists of four sequential stages: 
(1) audio preprocessing, 
(2) embedding extraction, 
(3) similarity indexing, and 
(4) ranking and evaluation.

Audio recordings are first standardized through a consistent preprocessing pipeline.
Embeddings are then extracted using models of varying capacity.
The resulting representations are indexed using a fixed similarity-search mechanism, after which ranked retrieval is performed and evaluated using the predefined ranking metrics.

All components other than embedding extraction are held constant across experimental conditions.
This design isolates embedding-model capacity as the primary experimental variable.


\section{Conditioning Assumptions}

All claims in this thesis are conditioned on specific contextual assumptions. 
The study is conducted under a small-scale dataset setting ($N$ limited), 
within a single-domain classical music archive, 
using retrieval-based proxy evaluation, 
and without incorporating user interaction modeling.

Accordingly, the findings should be interpreted within these constraints. 
These assumptions define the validity domain of the study, 
and results should not be generalized beyond similar small-scale, 
single-domain archival settings without further empirical validation.

\subsection*{Assumption of Proxy Validity}

This study assumes that metadata-derived labels (e.g., Composer and Character) 
serve as valid, though coarse, indicators of musical similarity within the archival context. 
While musical similarity is inherently continuous and multidimensional, 
these proxy labels provide a controlled and reproducible approximation suitable for systematic evaluation.



